\documentclass[a4paper,11pt]{article}

% ---- Encoding & fonts -------------------------------------------------------
\usepackage[T1]{fontenc}
\usepackage[utf8]{inputenc}
\usepackage{lmodern}
\usepackage{microtype}              % protrusion/kerning -> fewer overfull boxes

% ---- Geometry (~1in margins) ------------------------------------------------
\usepackage[a4paper,margin=1in,headheight=15pt,headsep=14pt]{geometry}

% ---- Math -------------------------------------------------------------------
\usepackage{amsmath,amssymb}

% ---- Graphics, tables, colour ----------------------------------------------
\usepackage{graphicx}
\usepackage{booktabs}
\usepackage[table,svgnames]{xcolor}

% ---- Callout boxes (breakable pill-tab tcolorboxes) -------------------------
\usepackage[most]{tcolorbox}
\tcbuselibrary{skins,breakable}

% ---- Callout glyphs ---------------------------------------------------------
\usepackage{pifont}
\IfFileExists{bbding.sty}{%
  \usepackage{bbding}%
  \newcommand{\glyphHand}{\Pointinghand}%
  \newcommand{\glyphPencil}{\Pencil}%
}{%
  \newcommand{\glyphHand}{\ding{43}}%
  \newcommand{\glyphPencil}{\ding{46}}%
}

% ---- Lists & spacing --------------------------------------------------------
\usepackage{enumitem}
\setlist{leftmargin=1.5em,topsep=4pt,itemsep=3pt}

\newlist{steps}{enumerate}{1}
\setlist[steps]{label=\textbf{Step \arabic*.},
  leftmargin=4.4em, labelwidth=3.8em, labelsep=0.6em, labelindent=0pt,
  align=left, topsep=5pt, itemsep=6pt}
\usepackage{parskip}                % blank-line paragraph spacing, no indent

% ---- Headers / footers ------------------------------------------------------
\usepackage{fancyhdr}

% ---- Section styling --------------------------------------------------------
\usepackage{titlesec}

% ---- Links (hidden) + URL line-breaking -------------------------------------
\usepackage[hidelinks]{hyperref}
\usepackage{xurl}
\hypersetup{breaklinks=true}

% =============================================================================
%  COLOURS — navy chrome + the FIVE taxonomy ramps
% =============================================================================
\definecolor{navy}{HTML}{21355E}        % banner + headings
\definecolor{navyrule}{HTML}{21355E}
\definecolor{inktext}{HTML}{1A202C}     % body ink
\definecolor{mutedtext}{HTML}{6B7280}   % header / meta grey

% Intuition — blue
\definecolor{intuFrame}{HTML}{2C5AA0}   \definecolor{intuFill}{HTML}{F4F7FC}
% Everyday — amber
\definecolor{everyFrame}{HTML}{8A5A1E}  \definecolor{everyFill}{HTML}{FBF1E3}
% Worked — green
\definecolor{workFrame}{HTML}{2E7D52}   \definecolor{workFill}{HTML}{F1F8F3}
% Watch out — red
\definecolor{watchFrame}{HTML}{B23A48}  \definecolor{watchFill}{HTML}{FBF1F2}
% Key takeaway — purple/violet
\definecolor{keyFrame}{HTML}{6A4C93}    \definecolor{keyFill}{HTML}{F4F0F9}
% Wait, what? — teal
\definecolor{waitFrame}{HTML}{0F766E}   \definecolor{waitFill}{HTML}{EEF7F5}

\color{inktext}
\hypersetup{linkcolor=navy,urlcolor=navy,citecolor=navy}

% =============================================================================
%  RUNNING HEADER / FOOTER
% =============================================================================
\newcommand{\CourseShort}{Deep Neural Networks}
\newcommand{\SessionNo}{8}
\newcommand{\LectureTitle}{Foundations of Neural Networks and CNNs}

\pagestyle{fancy}
\fancyhf{}
\fancyhead[L]{\footnotesize\color{mutedtext}\textsf{\CourseShort}}
\fancyhead[R]{\footnotesize\color{mutedtext}\textsf{Session \SessionNo\ \textperiodcentered\ \LectureTitle}}
\fancyfoot[L]{\footnotesize\color{mutedtext}\textsf{WILP, BITS Pilani}}
\fancyfoot[C]{\footnotesize\color{mutedtext}\thepage}
\renewcommand{\headrulewidth}{0.4pt}
\renewcommand{\footrulewidth}{0pt}
\renewcommand{\headrule}{\hbox to\headwidth{\color{navyrule!55}\leaders\hrule height \headrulewidth\hfill}}

% =============================================================================
%  SECTION HEADINGS
% =============================================================================
\titleformat{\section}
  {\sffamily\Large\bfseries\color{navy}}
  {\thesection}
  {0.8em}
  {}
  [{\vspace{2pt}\color{navy!70}\titlerule[0.8pt]}]
\titleformat{\subsection}
  {\sffamily\large\bfseries\color{navy!92}}
  {\thesubsection}{0.7em}{}
\titlespacing*{\section}{0pt}{16pt}{8pt}
\titlespacing*{\subsection}{0pt}{12pt}{4pt}

\newif\ifmlkfirstsec \mlkfirstsectrue
\newcommand{\sectionbreak}{\ifmlkfirstsec\mlkfirstsecfalse\else\clearpage\fi}

\newcommand{\secsub}[1]{%
  \noindent{\sffamily\bfseries\itshape\color{navy!85}\large #1}\par\medskip}

\newcommand{\flipanswer}[1]{\rotatebox[origin=c]{180}{\parbox{0.92\linewidth}{\small #1}}}

% =============================================================================
%  PILL-TAB CALLOUTS
% =============================================================================
\tcbset{
  housecallout/.style 2 args={
    enhanced, breakable,
    colback=#2, colframe=#1,
    boxrule=0.6pt, arc=2.4pt, outer arc=2.4pt,
    left=10pt, right=10pt, top=9pt, bottom=9pt,
    before skip=11pt, after skip=11pt,
    fontupper=\normalsize,
    coltitle=white,
    fonttitle=\bfseries\sffamily\small,
    attach boxed title to top left={xshift=6mm,yshift=-3mm},
    boxed title style={
      colback=#1, colframe=#1,
      rounded corners, arc=3pt,
      boxrule=0pt, left=6pt, right=6pt, top=2pt, bottom=2pt
    }
  }
}

\newtcolorbox{intuition}{%
  housecallout={intuFrame}{intuFill},
  title={\raisebox{-0.05ex}{\glyphHand}\,\ The intuition}}

\newtcolorbox{everyday}{%
  housecallout={everyFrame}{everyFill},
  title={\raisebox{-0.05ex}{$\bigstar$}\,\ Everyday picture}}

\newtcolorbox{worked}[1]{%
  housecallout={workFrame}{workFill},
  title={\raisebox{-0.05ex}{\glyphPencil}\,\ Worked example: #1}}

\newtcolorbox{watchout}{%
  housecallout={watchFrame}{watchFill},
  title={\raisebox{-0.05ex}{\ding{55}}\,\ Watch out}}

\newtcolorbox{keytake}{%
  housecallout={keyFrame}{keyFill},
  title={\raisebox{-0.05ex}{\ding{51}}\,\ Key takeaway}}

\newtcolorbox{waitwhat}{%
  housecallout={waitFrame}{waitFill},
  title={\raisebox{-0.05ex}{\ding{93}}\,\ Wait --- really?}}

\newif\ifbitslogo \bitslogofalse
\IfFileExists{assets/bits-logo.png}%
  {\bitslogotrue \def\bitslogopath{assets/bits-logo.png}}%
  {\IfFileExists{../../templates/assets/bits-logo.png}%
     {\bitslogotrue \def\bitslogopath{../../templates/assets/bits-logo.png}}%
     {}}
\newcommand{\bitswordmark}{{\sffamily\bfseries\color{white}\fontsize{12.5}{15}\selectfont WILP, BITS Pilani}}
\newcommand{\bitslogochip}{\ifbitslogo\colorbox{white}{\,\raisebox{-0.2ex}{\includegraphics[height=26pt]{\bitslogopath}}\,}\fi}

\providecommand{\addfontfeatureslike}[1]{#1}
\newcommand{\titlebanner}[4]{%
  \begin{tcolorbox}[
    enhanced, breakable=false,
    colback=navy, colframe=navy,
    boxrule=0pt, arc=3pt, outer arc=3pt,
    left=16pt, right=16pt, top=16pt, bottom=16pt,
    before skip=2pt, after skip=14pt,
    width=\linewidth ]
    \noindent
    \begin{minipage}[c]{0.72\linewidth}\bitswordmark\end{minipage}%
    \hfill
    \begin{minipage}[c]{0.26\linewidth}\raggedleft\bitslogochip\end{minipage}\par
    \vspace{9pt}
    {\sffamily\footnotesize\bfseries\color{white!85}%
      \MakeUppercase{\addfontfeatureslike{#1}}}\par
    \vspace{6pt}
    {\sffamily\bfseries\fontsize{30}{34}\selectfont\color{white} #2\par}
    \vspace{5pt}
    {\sffamily\large\color{white!88} #3\par}
    \vspace{9pt}
    {\color{white!55}\rule{\linewidth}{0.4pt}}\par
    \vspace{6pt}
    {\sffamily\footnotesize\color{white!82} #4\par}
  \end{tcolorbox}%
}

\newcommand{\housefig}[2]{%
  \begin{figure}[htbp]
    \centering
    \includegraphics[width=\linewidth]{#1}%
    \caption{#2}
  \end{figure}%
}

\newcommand{\vocab}[1]{\textbf{\color{navy}#1}}

\begin{document}

\thispagestyle{fancy}

\titlebanner
  {Deep Neural Networks \textperiodcentered\ Companion Reader}
  {From Perceptrons to Convolutional Architectures}
  {A complete guide to artificial neurons, linear models, multi-class classification, feed-forward networks, and CNN building blocks.}
  {Session 8 \textperiodcentered\ Read alongside the lecture slides \textperiodcentered\ Every slide example worked out in full}

\smallskip
\noindent{\small\color{mutedtext}%
Prepared for the Work Integrated Learning Programmes (WILP), BITS Pilani.}
\smallskip

\noindent\textbf{How to use this companion.}
The slides give you formulas and architecture diagrams. This companion explains the \emph{why} behind every equation and step. We build every model from simple weighted sums up to deep convolutional neural networks (CNNs). Colour-coded callout boxes guide your reading: \textcolor{intuFrame}{\textbf{blue}} for intuition, \textcolor{everyFrame}{\textbf{amber}} for real-life pictures, \textcolor{workFrame}{\textbf{green}} for fully worked numerical examples, \textcolor{watchFrame}{\textbf{red}} for common pitfalls, and \textcolor{keyFrame}{\textbf{purple}} for core takeaways.

\medskip

\section{First taste: The artificial neuron and the perceptron}
\secsub{Multiply features by weights, add a bias, and pass the sum through a threshold function.}

An \vocab{artificial neuron} is the atomic building block of deep learning. It takes an input feature vector $\mathbf{x} = [x_1, x_2, \dots, x_d]^T$, multiplies each feature by a weight $w_i$, adds a bias term $b$, and computes a weighted score $z$:
\begin{equation}
z = \sum_{i=1}^d w_i x_i + b = \mathbf{w}^T \mathbf{x} + b.
\end{equation}
The neuron then applies an \vocab{activation function} $f(z)$ to produce the final output $y = f(z)$.

The simplest form of an artificial neuron is the \vocab{perceptron}. It uses a step function as its activation function:
\begin{equation}
f(z) = \begin{cases} 1 & \text{if } z \ge 0 \\ 0 & \text{if } z < 0 \end{cases}
\end{equation}

\begin{intuition}
Think of a neuron as a committee member scoring a proposal. Each input $x_i$ is an evidence item. The weight $w_i$ is how much the member trusts that item. The bias $b$ is the baseline willingness to say yes. If the weighted total meets the threshold, the output fires.
\end{intuition}

\begin{everyday}
Deciding whether to attend an outdoor concert. Inputs: $x_1 = 1$ (good weather), $x_2 = 1$ (favourite band). If weather matters more ($w_1 = 3$) than the band ($w_2 = 2$) and your threshold is $b = -4$, you attend only if $3x_1 + 2x_2 - 4 \ge 0$.
\end{everyday}

A single perceptron can learn simple logical gates like \textbf{AND} or \textbf{OR}. For an AND gate with two inputs $x_1, x_2 \in \{0, 1\}$, set weights $w_1 = 1, w_2 = 1$ and bias $b = -1.5$. The weighted score is $z = 1\cdot x_1 + 1\cdot x_2 - 1.5$. When both inputs are $1$, $z = 0.5 \ge 0 \implies y = 1$. Otherwise, $z < 0 \implies y = 0$.

However, a single perceptron \textbf{fails on the XOR gate}. An XOR gate outputs $1$ when $x_1 \neq x_2$ and $0$ when $x_1 = x_2$.
Plotting the four input points $(0,0), (0,1), (1,0), (1,1)$ on a plane shows that no single straight line can separate the two classes. A single perceptron can only draw a linear decision boundary. To solve non-linearly separable problems like XOR, we must stack neurons into multi-layer networks.

\begin{worked}{Why a single perceptron cannot solve XOR}
Consider the XOR truth table points: $(0,0) \to 0$, $(1,1) \to 0$, $(0,1) \to 1$, $(1,0) \to 1$.
\begin{steps}
  \item \textbf{Write linear inequalities for output 0.}
        For $(0,0) \to 0$, we need $w_1(0) + w_2(0) + b < 0 \implies b < 0$.
        For $(1,1) \to 0$, we need $w_1 + w_2 + b < 0$.
  \item \textbf{Write linear inequalities for output 1.}
        For $(1,0) \to 1$, we need $w_1 + b \ge 0$.
        For $(0,1) \to 1$, we need $w_2 + b \ge 0$.
  \item \textbf{Sum the inequalities for output 1.}
        $(w_1 + b) + (w_2 + b) \ge 0 \implies w_1 + w_2 + 2b \ge 0$.
  \item \textbf{Derive the contradiction.}
        Since $b < 0$, we have $w_1 + w_2 + b > w_1 + w_2 + 2b \ge 0 \implies w_1 + w_2 + b > 0$.
        This directly contradicts $w_1 + w_2 + b < 0$ required for $(1,1)$.
        No weights $(w_1, w_2, b)$ exist.
\end{steps}
\end{worked}

\housefig{figures/fig_perceptron_xor.png}{Linear boundary for AND gate vs non-linear boundary requirement for XOR gate.}

\begin{watchout}
Do not confuse single perceptrons with multi-layer feed-forward networks. A single perceptron has no hidden layers and can only draw flat, linear decision lines.
\end{watchout}

\begin{keytake}
A single perceptron computes a linear decision boundary. It solves linearly separable logic like AND, but fails on non-linear patterns like XOR. Solving complex problems requires stacking hidden layers.
\end{keytake}


\section{Predicting numbers: Linear regression and optimization}
\secsub{Predict a continuous output as a weighted sum of inputs, and adjust weights using gradient descent.}

When our target $y \in \mathbb{R}$ is continuous (like house prices or temperature), we use \vocab{linear regression}. The mathematical model is:
\begin{equation}
\hat{y} = w_0 + w_1 x_1 + w_2 x_2 + \dots + w_d x_d = \mathbf{w}^T \mathbf{x} + b
\end{equation}
Linear regression can be viewed as a single artificial neuron with an \vocab{identity activation function} $f(z) = z$.

The identity function has key properties:
\begin{enumerate}[leftmargin=1.5em]
  \item Continuous output over $(-\infty, \infty)$.
  \item Differentiable everywhere with derivative $f'(z) = 1$.
  \item Allows gradients to flow backward unchanged during optimization.
\end{enumerate}

To measure how well our model fits the training data, we define an objective function called \vocab{squared error loss}:
\begin{equation}
L(y, \hat{y}) = \frac{1}{2} (y - \hat{y})^2
\end{equation}
For a dataset of $N$ examples, the Mean Squared Error (MSE) cost function is $J(\mathbf{w}, b) = \frac{1}{N} \sum_{i=1}^N \frac{1}{2} (y^{(i)} - \hat{y}^{(i)})^2$.

To minimize this loss, we compute the gradient with respect to weights and update them using Stochastic Gradient Descent (SGD):
\begin{align}
\frac{\partial L}{\partial w_j} &= \frac{\partial L}{\partial \hat{y}} \cdot \frac{\partial \hat{y}}{\partial z} \cdot \frac{\partial z}{\partial w_j} = - (y - \hat{y}) \cdot 1 \cdot x_j = (\hat{y} - y) x_j \\
w_j^{(\text{new})} &= w_j^{(\text{old})} - \eta \frac{\partial L}{\partial w_j}
\end{align}
where $\eta > 0$ is the learning rate.

To evaluate model performance, we use the \vocab{$R^2$ score} (Coefficient of Determination):
\begin{equation}
R^2 = 1 - \frac{SS_{\text{res}}}{SS_{\text{tot}}} = 1 - \frac{\sum_i (y_i - \hat{y}_i)^2}{\sum_i (y_i - \bar{y})^2}
\end{equation}
An $R^2 = 1$ indicates perfect prediction, while $R^2 = 0$ means the model performs no better than predicting the mean $\bar{y}$.

\begin{worked}{One SGD update step for Linear Regression}
Let a training sample be $\mathbf{x} = [2.0, 1.0]^T$ with true target $y = 5.0$. Initial weights are $\mathbf{w} = [0.5, -0.5]^T$, bias $b = 0.1$, and learning rate $\eta = 0.1$.
\begin{steps}
  \item \textbf{Compute prediction.}
        $\hat{y} = w_1 x_1 + w_2 x_2 + b = (0.5)(2.0) + (-0.5)(1.0) + 0.1 = 1.0 - 0.5 + 0.1 = 0.6$.
  \item \textbf{Compute initial loss.}
        $L = \frac{1}{2} (y - \hat{y})^2 = \frac{1}{2} (5.0 - 0.6)^2 = \frac{1}{2} (4.4)^2 = 9.68$.
  \item \textbf{Compute prediction error and gradients.}
        Error $(\hat{y} - y) = 0.6 - 5.0 = -4.4$. \\
        $\frac{\partial L}{\partial w_1} = (\hat{y} - y) x_1 = (-4.4)(2.0) = -8.8$. \\
        $\frac{\partial L}{\partial w_2} = (\hat{y} - y) x_2 = (-4.4)(1.0) = -4.4$. \\
        $\frac{\partial L}{\partial b} = (\hat{y} - y) = -4.4$.
  \item \textbf{Update weights and bias.}
        $w_1^{(\text{new})} = 0.5 - 0.1(-8.8) = 0.5 + 0.88 = 1.38$. \\
        $w_2^{(\text{new})} = -0.5 - 0.1(-4.4) = -0.5 + 0.44 = -0.06$. \\
        $b^{(\text{new})} = 0.1 - 0.1(-4.4) = 0.1 + 0.44 = 0.54$.
  \item \textbf{Compute new prediction and loss.}
        $\hat{y}_{\text{new}} = (1.38)(2.0) + (-0.06)(1.0) + 0.54 = 2.76 - 0.06 + 0.54 = 3.24$. \\
        $L_{\text{new}} = \frac{1}{2} (5.0 - 3.24)^2 = \frac{1}{2} (1.76)^2 = 1.5488$.
        Loss decreased from $9.68$ to $1.55$ in a single step!
\end{steps}
\end{worked}

\begin{keytake}
Linear regression predicts continuous targets using identity activation. We optimize weights by moving in the opposite direction of the loss gradient, reducing prediction error.
\end{keytake}


\section{Predicting categories: Logistic regression for binary classification}
\secsub{Squash weighted scores into probabilities using the sigmoid activation function.}

In \vocab{classification}, our goal is to assign inputs to discrete categories. For binary classification, the target is $y \in \{0, 1\}$.
Linear regression fails for classification because its outputs are unbounded $(-\infty, \infty)$ and sensitive to outliers.

\vocab{Logistic regression} solves this by passing the score $z = \mathbf{w}^T \mathbf{x} + b$ through the \vocab{sigmoid activation function} $\sigma(z)$:
\begin{equation}
\hat{y} = \sigma(z) = \frac{1}{1 + e^{-z}} = \frac{1}{1 + e^{-(\mathbf{w}^T \mathbf{x} + b)}}
\end{equation}
The sigmoid function maps any real number $z \in (-\infty, \infty)$ to a valid probability $\hat{y} \in (0, 1)$.

Properties of Sigmoid:
\begin{enumerate}[leftmargin=1.5em]
  \item Output ranges between $0$ and $1$.
  \item Smooth and differentiable everywhere, with derivative $\sigma'(z) = \sigma(z) (1 - \sigma(z))$.
  \item Thresholding: if $\hat{y} \ge 0.5 \implies \text{Class } 1$, else $\text{Class } 0$.
\end{enumerate}

To train logistic regression, we use the \vocab{Binary Cross-Entropy (BCE) loss}:
\begin{equation}
L(y, \hat{y}) = - \big[ y \log \hat{y} + (1 - y) \log(1 - \hat{y}) \big]
\end{equation}
When true label $y = 1$, loss is $-\log(\hat{y})$. If prediction $\hat{y} \to 1$, loss $\to 0$. If $\hat{y} \to 0$, loss $\to \infty$.

Differentiating BCE loss with respect to $z$ yields a remarkably clean gradient:
\begin{equation}
\frac{\partial L}{\partial z} = \hat{y} - y
\end{equation}
Hence, weight update for feature $x_j$ is simply:
\begin{equation}
w_j^{(\text{new})} = w_j - \eta (\hat{y} - y) x_j
\end{equation}

\housefig{figures/fig_logistic_sigmoid.png}{Sigmoid activation function mapping unbounded scores z into probability range (0, 1).}

To evaluate classification performance beyond simple accuracy, we use a \vocab{confusion matrix}:
\begin{center}
\small
\begin{tabular}{@{}ccc@{}}
\toprule
& \textbf{Predicted Positive (1)} & \textbf{Predicted Negative (0)} \\ \midrule
\textbf{Actual Positive (1)} & True Positive (TP) & False Negative (FN) \\
\textbf{Actual Negative (0)} & False Positive (FP) & True Negative (TN) \\ \bottomrule
\end{tabular}
\end{center}

Standard classification metrics are:
\begin{align}
\text{Accuracy} &= \frac{\text{TP} + \text{TN}}{\text{TP} + \text{TN} + \text{FP} + \text{FN}} \\
\text{Precision} &= \frac{\text{TP}}{\text{TP} + \text{FP}} \quad (\text{Of all predicted positives, how many were right?}) \\
\text{Recall} &= \frac{\text{TP}}{\text{TP} + \text{FN}} \quad (\text{Of all actual positives, how many did we catch?}) \\
F_1 \text{ Score} &= 2 \cdot \frac{\text{Precision} \cdot \text{Recall}}{\text{Precision} + \text{Recall}} \quad (\text{Harmonic mean of precision and recall})
\end{align}

\begin{worked}{Two SGD steps for Logistic Regression}
Given input $\mathbf{x} = [1.0, 2.0]^T$, true label $y = 1$, initial weights $\mathbf{w} = [0.1, -0.2]^T$, bias $b = 0.0$, and learning rate $\eta = 0.5$.
\begin{steps}
  \item \textbf{Iteration 1: Score and prediction.} \\
        $z^{(1)} = (0.1)(1.0) + (-0.2)(2.0) + 0.0 = 0.1 - 0.4 = -0.3$. \\
        $\hat{y}^{(1)} = \sigma(-0.3) = \frac{1}{1 + e^{0.3}} \approx \frac{1}{1 + 1.34986} \approx 0.4255$.
  \item \textbf{Iteration 1: Error and update.} \\
        Error $(\hat{y} - y) = 0.4255 - 1.0 = -0.5745$. \\
        $w_1^{(1)} = 0.1 - 0.5(-0.5745)(1.0) = 0.1 + 0.28725 = 0.38725$. \\
        $w_2^{(1)} = -0.2 - 0.5(-0.5745)(2.0) = -0.2 + 0.5745 = 0.3745$. \\
        $b^{(1)} = 0.0 - 0.5(-0.5745) = 0.28725$.
  \item \textbf{Iteration 2: Score and prediction with new weights.} \\
        $z^{(2)} = (0.38725)(1.0) + (0.3745)(2.0) + 0.28725 = 0.38725 + 0.7490 + 0.28725 = 1.4235$. \\
        $\hat{y}^{(2)} = \sigma(1.4235) = \frac{1}{1 + e^{-1.4235}} \approx \frac{1}{1 + 0.24087} \approx 0.8059$.
  \item \textbf{Check progress.} \\
        Predicted probability for class 1 increased from $0.4255$ to $0.8059$. The model is rapidly learning!
\end{steps}
\end{worked}

\begin{keytake}
Logistic regression uses sigmoid activation to output probabilities for binary classification. Binary cross-entropy loss penalizes confident wrong predictions.
\end{keytake}


\section{Multi-class classification and the Softmax function}
\secsub{Convert raw output scores for K classes into a normalized probability distribution.}

When a classification task has $K > 2$ discrete classes, we format target labels using \vocab{one-hot encoding}. For class $k$, label vector $\mathbf{y}$ has length $K$ with $1$ at position $k$ and $0$ elsewhere. For example, with $K=3$ classes: Class 1 is $[1,0,0]^T$, Class 2 is $[0,1,0]^T$, Class 3 is $[0,0,1]^T$.

For $K$ classes, the network output layer produces $K$ raw real-valued scores called \vocab{logits} $\mathbf{z} = [z_1, z_2, \dots, z_K]^T$. We convert logits into class probabilities $P(y=k \mid \mathbf{x})$ using the \vocab{Softmax function}:
\begin{equation}
\hat{y}_k = \text{Softmax}(z_k) = \frac{e^{z_k}}{\sum_{j=1}^K e^{z_j}}
\end{equation}

Properties of Softmax:
\begin{enumerate}[leftmargin=1.5em]
  \item Every output $\hat{y}_k > 0$.
  \item The sum of all class probabilities equals 1: $\sum_{k=1}^K \hat{y}_k = 1$.
  \item Exponentials amplify differences between top logits, making the dominant class stand out.
\end{enumerate}

The loss function for multi-class classification is \vocab{Categorical Cross-Entropy Loss}:
\begin{equation}
L(\mathbf{y}, \hat{\mathbf{y}}) = - \sum_{k=1}^K y_k \log \hat{y}_k = - \log \hat{y}_c
\end{equation}
where $c$ is the true class label.

\housefig{figures/fig_softmax.png}{Softmax operation converting raw logits z into a sum-to-one probability distribution.}

\begin{worked}{Softmax computation and cross-entropy loss}
Given raw logits $\mathbf{z} = [2.0, 1.0, 0.1]^T$ for a 3-class problem where true class is Class 1 ($y = [1, 0, 0]^T$).
\begin{steps}
  \item \textbf{Compute exponentials.} \\
        $e^{z_1} = e^{2.0} \approx 7.3891$. \\
        $e^{z_2} = e^{1.0} \approx 2.7183$. \\
        $e^{z_3} = e^{0.1} \approx 1.1052$.
  \item \textbf{Compute sum of exponentials.} \\
        $S = 7.3891 + 2.7183 + 1.1052 = 11.2126$.
  \item \textbf{Compute Softmax probabilities.} \\
        $\hat{y}_1 = 7.3891 / 11.2126 \approx 0.6590 \quad (65.9\%)$. \\
        $\hat{y}_2 = 2.7183 / 11.2126 \approx 0.2424 \quad (24.2\%)$. \\
        $\hat{y}_3 = 1.1052 / 11.2126 \approx 0.0986 \quad (9.9\%)$. \\
        Check sum: $0.6590 + 0.2424 + 0.0986 = 1.0000$.
  \item \textbf{Compute Categorical Cross-Entropy Loss.} \\
        $L = - \log(\hat{y}_1) = - \log(0.6590) \approx 0.4170$.
\end{steps}
\end{worked}

\begin{keytake}
Softmax converts unbounded logits into a normalized probability distribution across multiple classes. Combined with cross-entropy loss, it drives model training for multi-class classification.
\end{keytake}


\section{Deep Feed-Forward Neural Networks (DFNN)}
\secsub{Stack hidden layers of non-linear transformations to build expressive deep representations.}

A \vocab{Deep Feed-Forward Neural Network (DFNN)} consists of an input layer, one or more hidden layers, and an output layer. Information moves strictly forward through the network.

For layer $l \in \{1, 2, \dots, L\}$:
\begin{align}
\mathbf{z}^{(l)} &= \mathbf{W}^{(l)} \mathbf{a}^{(l-1)} + \mathbf{b}^{(l)} \\
\mathbf{a}^{(l)} &= f^{(l)}(\mathbf{z}^{(l)})
\end{align}
where $\mathbf{a}^{(0)} = \mathbf{x}$ is the input feature vector, $\mathbf{W}^{(l)}$ is the weight matrix of shape $(n_l \times n_{l-1})$, $\mathbf{b}^{(l)}$ is the bias vector of length $n_l$, and $f^{(l)}$ is a non-linear activation function (like ReLU or Sigmoid).

Calculating total learnable parameters in a feed-forward network:
For layer $l$ connected to layer $l-1$:
\begin{equation}
\text{Params}^{(l)} = n_l \times n_{l-1} + n_l = n_l (n_{l-1} + 1)
\end{equation}
Total network parameters equal the sum of parameters across all layers.

\begin{worked}{Forward propagation through a 2-layer network}
Let input $\mathbf{x} = [1, 2]^T$. \\
Layer 1 (Hidden, 2 units, ReLU activation): $\mathbf{W}^{(1)} = \begin{bmatrix} 1 & 0 \\ -1 & 2 \end{bmatrix}$, $\mathbf{b}^{(1)} = \begin{bmatrix} 0 \\ 1 \end{bmatrix}$. \\
Layer 2 (Output, 1 unit, Linear activation): $\mathbf{W}^{(2)} = \begin{bmatrix} 2 & -1 \end{bmatrix}$, $b^{(2)} = 0.5$.
\begin{steps}
  \item \textbf{Hidden layer weighted sum.} \\
        $\mathbf{z}^{(1)} = \mathbf{W}^{(1)} \mathbf{x} + \mathbf{b}^{(1)} = \begin{bmatrix} 1(1) + 0(2) \\ -1(1) + 2(2) \end{bmatrix} + \begin{bmatrix} 0 \\ 1 \end{bmatrix} = \begin{bmatrix} 1 \\ 3 \end{bmatrix} + \begin{bmatrix} 0 \\ 1 \end{bmatrix} = \begin{bmatrix} 1 \\ 4 \end{bmatrix}$.
  \item \textbf{Hidden layer activation (ReLU).} \\
        $\mathbf{a}^{(1)} = \text{ReLU}\left(\begin{bmatrix} 1 \\ 4 \end{bmatrix}\right) = \begin{bmatrix} 1 \\ 4 \end{bmatrix}$.
  \item \textbf{Output layer weighted sum.} \\
        $z^{(2)} = \mathbf{W}^{(2)} \mathbf{a}^{(1)} + b^{(2)} = \begin{bmatrix} 2 & -1 \end{bmatrix} \begin{bmatrix} 1 \\ 4 \end{bmatrix} + 0.5 = (2 - 4) + 0.5 = -1.5$.
  \item \textbf{Final network prediction.} \\
        Since output activation is identity, $\hat{y} = a^{(2)} = -1.5$.
\end{steps}
\end{worked}

\noindent\textbf{Why depth matters.}
Deep networks learn \vocab{hierarchical feature representations}: early layers learn simple low-level primitives (edges, color blobs), middle layers combine primitives into shapes and parts, and deep layers assemble parts into complex high-level concepts (faces, cars).
Theoretical results show that deep networks can represent complex functions exponentially more efficiently than shallow, wide networks with the same number of parameters.

\noindent\textbf{Challenges in training deep networks:}
\begin{enumerate}[leftmargin=1.5em]
  \item \textbf{Vanishing Gradients:} Gradients shrink exponentially as they backpropagate through deep layers, causing early layers to stop updating. Solution: ReLU activations, residual skip connections, proper weight initialization.
  \item \textbf{Exploding Gradients:} Gradients blow up exponentially, causing numerical instability ($NaN$ values). Solution: Gradient clipping, normalization layers.
  \item \textbf{Degradation Problem:} Beyond a certain depth, standard deep networks suffer higher training error. Solution: Residual connections (ResNets).
\end{enumerate}

\begin{keytake}
Deep networks compose simple linear transformations and non-linear activations to learn rich hierarchical representations. Residual connections enable training networks over 100 layers deep.
\end{keytake}


\section{Convolutional Neural Networks (CNNs) and parameter efficiency}
\secsub{Preserve spatial structure using localized filters, parameter sharing, and spatial downsampling.}

Standard fully connected networks fail on high-dimensional grid inputs like images (e.g., $224 \times 224 \times 3$ image flattened into a $150,528$-dimensional vector) because parameter counts explode, causing severe overfitting and losing 2D spatial relationships.

\vocab{Convolutional Neural Networks (CNNs)} solve this using three spatial design principles:
\begin{enumerate}[leftmargin=1.5em]
  \item \textbf{Local Receptive Fields:} Filters look at small local patches of the input image at a time.
  \item \textbf{Parameter Sharing:} The same filter weights slide across the entire input image.
  \item \textbf{Spatial Invariance:} Features detected in one corner can be recognized anywhere in the image.
\end{enumerate}

\noindent\textbf{Convolution Output Spatial Dimensions Formula:}
For an input grid of height/width $N$, kernel size $K$, padding $P$, and stride $S$:
\begin{equation}
\text{Output Shape } O = \left\lfloor \frac{N + 2P - K}{S} \right\rfloor + 1
\end{equation}

\noindent\textbf{Types of Convolution Operations:}
\begin{itemize}[leftmargin=1.5em]
  \item \textbf{Standard Convolution:} Sliding a $K \times K \times C_{\text{in}}$ filter across spatial dimensions.
  \item \textbf{Dilated (Atrous) Convolution:} Inserts spaces (dilation rate $r$) between filter elements to expand receptive field without increasing parameter count. Effective kernel size becomes $K' = K + (K - 1)(r - 1)$.
  \item \textbf{Multi-Channel Convolution:} Input has shape $(H \times W \times C_{\text{in}})$. Each of the $C_{\text{out}}$ filters has shape $(K \times K \times C_{\text{in}})$. The filter convolves across all input channels simultaneously and sums the result to produce one 2D output feature map per filter.
\end{itemize}

\noindent\textbf{Pooling Operations (Spatial Downsampling):}
Pooling reduces spatial dimensions $(H, W)$ while preserving feature depth $C$.
\begin{itemize}[leftmargin=1.5em]
  \item \textbf{Max Pooling:} Extracts the maximum value in each window. Preserves prominent features and provides local translation invariance.
  \item \textbf{Average Pooling:} Computes average value in each window. Smooths downsampling.
  \item \textbf{Crucial Property:} Pooling layers have \textbf{ZERO learnable parameters}!
\end{itemize}

\housefig{figures/fig_conv_pool.png}{Complete pipeline: 6x6 image convolved with 3x3 filter, passed through ReLU, and max-pooled.}

\begin{worked}{Convolution, ReLU, and Max Pooling step-by-step}
Given a $6 \times 6$ binary input matrix $I$, a $3 \times 3$ filter $W$ with stride $S=1, P=0$, followed by ReLU and $2 \times 2$ Max Pooling with stride $S=2$:
$$
I = \begin{bmatrix}
1 & 0 & 0 & 0 & 0 & 1 \\
0 & 1 & 0 & 0 & 1 & 0 \\
0 & 0 & 1 & 1 & 0 & 0 \\
1 & 0 & 0 & 0 & 1 & 0 \\
0 & 1 & 0 & 0 & 1 & 0 \\
0 & 0 & 1 & 0 & 1 & 0
\end{bmatrix}, \quad
W = \begin{bmatrix}
-1 & 1 & -1 \\
-1 & 1 & -1 \\
-1 & 1 & -1
\end{bmatrix}
$$
\begin{steps}
  \item \textbf{Calculate output shape after convolution.} \\
        $O = \frac{6 - 3}{1} + 1 = 4 \implies 4 \times 4$ feature map.
  \item \textbf{Compute Convolution at top-left patch $I[0:3, 0:3]$.} \\
        Patch $= \begin{bmatrix} 1 & 0 & 0 \\ 0 & 1 & 0 \\ 0 & 0 & 1 \end{bmatrix}$. \\
        Dot product: $(1)(-1) + (0)(1) + (0)(-1) + (0)(-1) + (1)(1) + (0)(-1) + (0)(-1) + (0)(1) + (1)(-1) = -1 + 1 - 1 = -1$.
  \item \textbf{Convolved $4 \times 4$ Feature Map Result.} \\
        $Z = \begin{bmatrix} -1 & -1 & -1 & -1 \\ -1 & -1 & -2 & -1 \\ -1 & -1 & -2 & 1 \\ -1 & 0 & -4 & 3 \end{bmatrix}$.
  \item \textbf{Apply ReLU Activation $\max(0, z)$.} \\
        $A = \text{ReLU}(Z) = \begin{bmatrix} 0 & 0 & 0 & 0 \\ 0 & 0 & 0 & 0 \\ 0 & 0 & 0 & 1 \\ 0 & 0 & 0 & 3 \end{bmatrix}$.
  \item \textbf{Apply $2 \times 2$ Max Pooling with Stride 2.} \\
        Output shape: $O_{\text{pool}} = \frac{4 - 2}{2} + 1 = 2 \implies 2 \times 2$. \\
        Top-left block $\begin{bmatrix} 0 & 0 \\ 0 & 0 \end{bmatrix} \to \max = 0$. \\
        Bottom-right block $\begin{bmatrix} 0 & 1 \\ 0 & 3 \end{bmatrix} \to \max = 3$. \\
        Final Pooled Output: $\begin{bmatrix} 0 & 0 \\ 0 & 3 \end{bmatrix}$.
\end{steps}
\end{worked}


\section{Parameter counting and CNN architecture design patterns}
\secsub{Master exact parameter calculations and architectural patterns for high-performance networks.}

Understanding parameter counting is critical for designing efficient neural networks.

\noindent\textbf{1. Convolutional Layer Parameter Formula:}
For a Conv layer with kernel size $K \times K$, $C_{\text{in}}$ input channels, $C_{\text{out}}$ output filters, with bias:
\begin{equation}
\text{Params}_{\text{Conv}} = \big( K \times K \times C_{\text{in}} + 1 \big) \times C_{\text{out}}
\end{equation}
Note: Parameters depend ONLY on kernel dimensions and channels, \textbf{NOT} on spatial height or width!

\noindent\textbf{2. Pooling Layer Parameter Formula:}
\begin{equation}
\text{Params}_{\text{Pool}} = 0
\end{equation}

\noindent\textbf{3. Fully Connected (FC) Layer Parameter Formula:}
For an input vector of length $n$ connected to $m$ output neurons:
\begin{equation}
\text{Params}_{\text{FC}} = m \times (n + 1) = m \cdot n + m
\end{equation}

\begin{worked}{Full CNN Parameter Calculation}
Consider a CNN with:
Input: $32 \times 32 \times 3$ image. \\
Conv1: 64 filters of size $5 \times 5$, stride 1, padding 2. \\
Pool1: $2 \times 2$ max pool, stride 2. \\
Conv2: 128 filters of size $3 \times 3$, stride 1, padding 1. \\
FC1: Dense layer with 256 units. \\
Output FC: 10 units.
\begin{steps}
  \item \textbf{Conv1 parameters and output shape.} \\
        $\text{Params}_{\text{Conv1}} = (5 \times 5 \times 3 + 1) \times 64 = (75 + 1) \times 64 = 76 \times 64 = 4,864$. \\
        Output shape: $\left(\frac{32 + 2(2) - 5}{1} + 1\right) = 32 \implies 32 \times 32 \times 64$.
  \item \textbf{Pool1 parameters and output shape.} \\
        $\text{Params}_{\text{Pool1}} = 0$. \\
        Output shape: $16 \times 16 \times 64$.
  \item \textbf{Conv2 parameters and output shape.} \\
        $\text{Params}_{\text{Conv2}} = (3 \times 3 \times 64 + 1) \times 128 = (576 + 1) \times 128 = 577 \times 128 = 73,856$. \\
        Output shape: $16 \times 16 \times 128$.
  \item \textbf{Flatten step before FC1.} \\
        Flatten $16 \times 16 \times 128 = 32,768$ input dimensions.
  \item \textbf{FC1 parameters.} \\
        $\text{Params}_{\text{FC1}} = (32,768 + 1) \times 256 = 8,388,864$.
  \item \textbf{Output FC parameters.} \\
        $\text{Params}_{\text{Output}} = (256 + 1) \times 10 = 2,570$.
  \item \textbf{Total network parameters.} \\
        $4,864 + 0 + 73,856 + 8,388,864 + 2,570 = 8,470,154$ parameters ($\approx 8.47 \text{ M}$).
\end{steps}
\end{worked}

\noindent\textbf{Key Architectural Design Patterns in Modern CNNs:}
\begin{enumerate}[leftmargin=1.5em]
  \item \textbf{Spatial Reduction + Channel Expansion:} As spatial dimensions $(H, W)$ decrease via pooling, channel depth $C$ increases (e.g., $32\times32\times64 \to 16\times16\times128 \to 8\times8\times256$). This maintains feature capacity.
  \item \textbf{Repetitive Building Blocks:} Stacking standard modular blocks (like VGG 3x3 Conv blocks) simplifies network design.
  \item \textbf{Skip/Residual Connections (ResNet):} Adds identity shortcut connections $y = F(\mathbf{x}) + \mathbf{x}$. Enables backpropagation through 100+ layers by letting gradients flow directly through identity paths.
  \item \textbf{Bottleneck Design:} Uses $1 \times 1$ convolutions to shrink channel dimensions before expensive $3 \times 3$ convolutions, reducing parameters by up to $8\times$ without losing expressiveness.
  \item \textbf{Multi-Scale Processing (Inception):} Applies $1\times1, 3\times3, 5\times5$ filters in parallel and concatenates outputs to capture features at multiple receptive field scales.
\end{enumerate}

\begin{keytake}
Convolutional layers achieve immense parameter efficiency via weight sharing. Modern CNN architectures rely on spatial reduction, residual skip connections, and bottleneck designs to train deep models stably.
\end{keytake}


\section*{Glossary \& symbols}
\addcontentsline{toc}{section}{Glossary \& symbols}

\noindent\textbf{Key terms.}
\begin{description}[leftmargin=9.5em,style=nextline,font=\normalfont\bfseries\color{navy}]
  \item[Artificial Neuron] Basic computational unit performing weighted sum followed by non-linear activation.
  \item[Perceptron] Simplest neuron with step function activation; computes linear decision boundaries.
  \item[Sigmoid] S-shaped activation function mapping scores into range $(0,1)$ for probabilities.
  \item[Softmax] Function converting unnormalized logits into normalized multi-class probability distribution.
  \item[Binary Cross-Entropy] Loss function measuring divergence between predicted probabilities and binary labels.
  \item[Convolution Layer] Layer applying sliding learnable filters to preserve spatial feature structure.
  \item[Max Pooling] Fixed operation extracting maximum value in spatial windows with zero parameters.
  \item[Residual Connection] Identity shortcut connection $y = F(\mathbf{x}) + \mathbf{x}$ preventing vanishing gradients.
\end{description}

\medskip
\noindent\textbf{Symbols at a glance.}
\begin{center}\small
\begin{tabular}{@{}ll@{}}
\toprule \textbf{Symbol} & \textbf{Meaning} \\ \midrule
$\mathbf{w}, b$ & Weights vector and scalar bias term \\
$\sigma(z)$ & Sigmoid activation function $\frac{1}{1 + e^{-z}}$ \\
$L(y, \hat{y})$ & Loss function measuring discrepancy between true label $y$ and prediction $\hat{y}$ \\
$\eta$ & Learning rate parameter for gradient descent updates \\
$K, P, S$ & Convolution kernel size, padding width, and stride step size \\
\bottomrule
\end{tabular}
\end{center}

\end{document}
